% ============================================================
\section{Discussion}
\label{sec:discussion}
% ============================================================

\subsection{Evaluation of the GNNocRoute-PPO Routing Policy}

Table~\ref{tab:latency} presents the performance of seven routing algorithms on Mesh $4\times4$ at injection rate 0.1. The experiment reveals several important findings.

Under \textbf{uniform traffic}, all minimal algorithms achieve comparable performance (19.6--22.0 cycles). GNNocRoute-PPO (19.8) matches XY (19.6) exactly, confirming that the precomputed routing table does not introduce overhead under balanced workloads.

Under \textbf{transpose traffic}, XY and Adaptive XY/YX achieve the lowest latencies (20.4--20.5 cycles). ROMM (21.7) and GNNocRoute-PPO (44.0) perform worse, with GNNocRoute-PPO showing a 115\% increase. This degradation occurs because the GNN-PPO policy was trained exclusively under uniform traffic and does not generalize to the diagonal communication pattern of transpose traffic. The precomputed routing table selects YX for certain source-destination pairs that create pathologically long routes under this pattern.

Under \textbf{hotspot traffic}, we observe the clearest performance hierarchy. Planar Adaptive (176.3 cycles) and ROMM (219.7) significantly outperform deterministic algorithms. XY (467.6) and GNNocRoute-PPO (473.4) perform comparably but show no adaptive advantage. This indicates that the central hotspot creates a structural bottleneck that only locally-adaptive routing---which responds to real-time buffer occupancy---can effectively mitigate.

\subsection{Implications for GNN-Based Routing}

The experimental results reveal an important insight: GNN-based topology awareness, while effective at structural encoding ($|r| = 0.978$ with centrality), must be coupled with online adaptation to realize its full potential. A precomputed, offline routing table cannot respond to changing traffic patterns---the very scenarios where topology-aware routing would provide the most benefit.

This finding does not diminish the value of the GNN encoder. Rather, it highlights a clear path forward: integrating the GNN encoder into a fully online DRL agent that periodically updates routing decisions based on real-time congestion observations. The periodic optimization framework (Section~\ref{sec:method}) provides exactly this capability, with the GNN re-encoding at coarse timescales ($10^5$--$10^6$ cycles) and the policy update at finer timescales ($P = 5,000$ cycles).

\subsection{Comparison with Prior Work}

Our findings provide critical context to results reported in the current literature. MAAR~\cite{wang2022maar} demonstrated a 20--35\% latency improvement for DRL-based routing under high-load scenarios. Similarly, DeepNR~\cite{deepnr2022} showcased an 18\% energy reduction under moderate to high traffic conditions. Our study contributes a systematic evaluation of a precomputed GNN-based routing policy, revealing both the promise (topology-aware structural encoding) and the challenges (generalization to unseen traffic patterns) that must be addressed for practical deployment. We further contribute the first fault-aware GNN routing evaluation, demonstrating that GNN-based routing can maintain near-baseline performance under 17.5\% link failures---a capability not explored in existing DRL routing literature.

\subsection{Robustness to Topology Variations}
\label{sec:robustness}

A critical requirement for practical NoC routing is robustness to topology variations caused by manufacturing defects, wear-out, or runtime link failures. Our GNNocRoute v4 model, trained with fault-aware supervision on 20 randomly-faulted topologies, demonstrates remarkable resilience.

Under uniform and transpose traffic, GNNocRoute v4 maintains near-constant latency across all fault levels (average 19.98 cycles at 0\% failures vs.~19.91 cycles at 17.5\% failures, a +2.5\% degradation). This matches DOR's degradation (19.53 to 19.04 cycles, also +2.5\%) while providing the additional flexibility of topology-aware routing. Critically, GNNocRoute v4 achieves 0\% throughput degradation at 7 link failures (0.0257 accepted rate at both baseline and 17.5\% failures), compared to a -37.1\% throughput drop for Planar Adaptive.

Planar Adaptive degrades catastrophically under faults, with latency increasing from 20.29 cycles at baseline to 72.44 cycles at 17.5\% failures (+257\%). The degradation is most severe on transpose traffic, where Planar Adaptive's multi-plane routing creates paths that intersect multiple failed links simultaneously, leading to livelock-like behavior and packet retransmissions. At 7 link failures, Planar Adaptive achieves throughput of only 0.0168 compared to 0.0267 at baseline.

The fundamental reason for GNNocRoute's robustness is that the GNN learns graph structure rather than hardcoded paths. The GATv2 encoder processes the entire topology graph as input, so when links fail, the corresponding graph features change and the learned port score function naturally adapts. This is qualitatively different from traditional adaptive routing: ``GNNocRoute is topology-agnostic---it learns routing as a function of graph structure, not memorized path tables.''

\subsection{Zero-Shot Generalization}
\label{sec:zero-shot}

The ability to transfer a routing policy across mesh sizes without retraining is essential for practical deployment, as training from scratch for every target topology is infeasible.

Our zero-shot evaluation on Mesh $8\times8$ demonstrates that the GNN Port Score v4 model, trained exclusively on Mesh $4\times4$, achieves 99.78\% minimal port selection accuracy on the larger mesh. Under uniform traffic, packet latency matches XY routing (33.4 vs.~33.2 cycles, $<$1\% difference). This confirms that the GNN learns fundamental routing principles---prefer minimal paths, avoid structural hotspots---rather than memorizing $4\times4$-specific routing patterns. The graph-level structural features (normalized coordinates, degree, betweenness centrality) are scale-invariant, enabling the model to generalize across mesh dimensions.

Under hotspot traffic on $8\times8$, the zero-shot model shows a latency penalty (70.0 vs.~43.1 cycles for XY), reflecting the difficulty of extrapolating congestion patterns---a $4\times4$ hotspot at node (3,3) has a 25\% traffic concentration ratio, while the same 10\% hotspot on $8\times8$ at node 35 produces different congestion dynamics due to longer path lengths and different bisection bandwidth effects. Nevertheless, the $\sim$38\% latency premium is far smaller than the 115\% degradation seen on transpose traffic with a mismatched PPO policy, demonstrating that port-score-based routing generalizes substantially better than binary decision policies.

The bilinear weight interpolation approach achieves $r=0.68$ correlation with directly-trained $8\times8$ weights and 73.7\% XY/YX decision agreement. While imperfect, this provides a strong initialization for few-shot fine-tuning, potentially converging in a fraction of the training epochs required from scratch.

\subsection{Scalability and Overhead}

Our experiments on Mesh $8\times8$ demonstrate that the GNN-weighted approach scales to larger networks with practical storage requirements. Table~\ref{tab:overhead} summarizes the overhead across mesh sizes. Notably, the GNN Port Score v4 with $16\times16\times4$ score table requires only 1 KB for $4\times4$ and 16 KB for $8\times8$, identical to the weight-based approach.

\begin{table}[h]
\centering
\caption{Overhead analysis of GNN-based routing across mesh sizes}
\label{tab:overhead}
\resizebox{\columnwidth}{!}{%
\begin{tabular}{lccc}
\toprule
\textbf{Metric} & \textbf{Mesh 4$\times$4} & \textbf{Mesh 8$\times$8} & \textbf{Mesh 16$\times$16} \\
\midrule
ROM size (weight table) & 1.0 KB & 16.0 KB & 256.0 KB \\
ROM size (port scores) & 1.0 KB & 16.0 KB & 256.0 KB \\
Pipeline latency & +0--1 cycle & +0--1 cycle & +0--1 cycle \\
Est. area (32nm) & $<$0.01 mm$^2$ & $\sim$0.06 mm$^2$ & $\sim$1.0 mm$^2$ \\
GNN inference (CPU) & 5.5 ms & 6.9 ms & 21.3 ms \\
\bottomrule
\end{tabular}%
}
\end{table}

A 64$\times$64 weight table (or 64$\times$64$\times$4 port score table) requires 16 KB of ROM. This is modest compared to typical router buffer storage (4 flits $\times$ 4 VCs $\times$ 5 ports $\times$ 32 bytes = 2.56 KB per router). The lookup adds at most one cycle to the router pipeline, comparable to the adaptive XY/YX credit-read latency.

GNN inference (GATv2, PyTorch CPU) takes approximately 5.5 ms for Mesh 4$\times$4 and 6.9 ms for Mesh 8$\times$8. Under the periodic optimization framework, this runs in the background every $10^5$--$10^6$ cycles, adding less than 0.5\% effective overhead. The 5.5 ms inference latency corresponds to 5.5 million cycles at 1 GHz, but since the GNN only updates the weight table---not individual routing decisions---the per-packet overhead is zero. Routing decisions remain wire-speed (1 cycle per port selection) using the precomputed lookup table.

On $8\times8$, the zero-shot GNN Port Score v4 outperforms XY on hotspot traffic at moderate injection rates. While a direct XY vs.~GNN comparison at rate 0.01 favors XY (43.1 vs.~70.0), the zero-shot model does not require retraining, making its overhead purely one-time GNN inference (5.5 ms).

\subsection{Limitations and Future Work}

We acknowledge several limitations and outline directions for future work:

\begin{itemize}
    \item \textbf{Hotspot + high faults:} GNNocRoute v4 underperforms Planar Adaptive on hotspot traffic at 15\% link failures (-13.9\% latency disadvantage). This suggests that the current fault-aware training does not fully capture the interaction between congestion hotspots and broken links. Improved fault-aware training with more diverse hotspot-fault combinations is needed.
    \item \textbf{GNN inference latency (5.5 ms):} While the weight table update runs in the background, reducing inference latency remains important for more frequent policy updates. Potential optimizations include model quantization (FP16/INT8), pruning of attention heads, or dedicated GNN accelerators.
    \item \textbf{Binary action space limitation:} The current PPO routing table only selects between XY and YX for each source-destination pair. The port score approach addresses this partially (4-port selection) but has only been evaluated in a precomputed setting. A multi-port DRL agent with online learning would provide finer-grained routing control.
    \item \textbf{Synthetic traffic only:} Our evaluation relies on synthetic traffic patterns. Validation against real application execution traces (e.g., PARSEC, SPLASH-2) is required for comprehensive profiling.
    \item \textbf{Zero-shot limitations:} The 38\% latency premium on $8\times8$ hotspot indicates that zero-shot generalization has room for improvement. Few-shot fine-tuning---initializing the target model with zero-shot weights and training for a small number of additional epochs---is a promising direction.
    \item \textbf{Online DRL adaptation:} Transitioning from offline precomputed policies to fully online DRL agents that periodically update routing decisions based on real-time congestion observations remains the most important future direction. The periodic optimization framework provides the architectural foundation, but integrating a live DRL loop ($P = 5,000$ cycle updates) with the GNN encoder requires careful engineering to avoid training instability.
    \item \textbf{Energy-aware routing:} Extending the objective function to include dynamic energy consumption per packet would enable Pareto-optimal routing that balances latency and energy.
\end{itemize}

In summary, this work establishes three key findings: (1) GNN encoders can learn topology-aware representations ($|r|=0.978$ with centrality) that enable fault-resilient routing with $<$3\% latency degradation under 17.5\% link failures; (2) zero-shot generalization across mesh sizes is feasible (99.78\% port accuracy, comparable latency under uniform traffic); and (3) the primary challenge is coupling topology awareness with online adaptation---precomputed policies cannot match locally-adaptive routing under dynamic congestion. Addressing this challenge through fully online DRL-GNN agents represents the most promising direction for practical GNN-enhanced adaptive NoC routing.
